
Toda reacción química toma un tiempo dado en ocurrir, sin embargo, existen sustancias capaces de acelerar la velocidad de estas reacciones, estas sustancias son conocidas como \emph{catalizadores}. Para la vida en la Tierra este tema de la velocidad de las reacciones químicas es de suma importancia. Tomemos como ejemplo, la degradación del peróxido de hidrógeno (\ce{H2O2}) en agua y oxígeno: 

\begin{equation*}
    \ce{H2O2 -> 2H2O + O2}
\end{equation*}

Esta reacción puede ocurrir de manera autónoma, sin ayuda externa de energía o sustancias extra, sin embargo, tomaría un tiempo considerable en ocurrir de esta forma. Dado que el \ce{H2O2} es un compuesto oxidante, puede ser dañino para los seres vivos, es necesario acelerar su reacción de degradación para deshacerse de él lo más rápido posible, o sea, necesitan un catalizador. Afortunadamente los seres vivos han desarrollado sus propios catalizadores, las \emph{enzimas}. Por ejemplo, el uso de la enzima como catalasa puede hacer que la degradación de peróxido mencionada antes ocurra hasta 1000 millones de veces más rápido\citep{mathewsBioquimica2002}, y así como esta enzima para esta reacción concreta, los seres vivos utilizan estas enzimas para hacer que sus procesos bioquímicos se desarrollen de manera adecuada.

Así pues, las enzimas se pueden definir como \emph{catalizadores biológicos}, moléculas que utilizan los seres vivos para acelerar reacciones químicas.

Más específicamente hablando, las enzimas son \emph{proteínas}. Conviene pues recordar la estructura básica de las proteínas, pues muchos de los aspectos que se estudian acerca de las enzimas, así como los procesos que se siguen para su aislamiento y uso, derivan de las características bioquímicas de las proteínas en general.



\section{Recordando la estructura de las proteínas}\label{sec:estructura_proteica_de_las_enzimas}

Las proteínas son \emph{polímeros}, o sea, sustancias compuestas de pequeñas moléculas ``base'' que se ensamblan entre sí. las cuales son llamados \emph{monómeros}. En el caso de las proteínas, sus monómeros son los \emph{aminoácidos}. Aspectos como la cantidad, el tipo y el orden de los aminoácidos de una proteína definen las características que esta tendrá.

Podemos dividir la \emph{estructura general} de un aminoácido en un \emph{carbono central alfa} ($\alpha$) del cual se desprenden 3 grupos: el grupo amino (\ce{NH2}), el grupo carboxilo (\ce{COOH}) y la cadena lateral (R). 

Los \emph{grupos amino y carboxilo} son grupos funcionales\footnote{O sea, capaces de reaccionar} a través de los cuales los aminoácidos se pueden unir entre sí. Tal unión entre aminoácidos por medio de tales grupos se llama \emph{enlace peptídico}. Estos grupos dotan además a los aminoácidos de un \emph{caracter anfotérico}, o sea, que puede actuar tanto como ácido y como base, esto porque cada grupo se ioniza de manera diferente dependiendo el pH del medio. 

Por otro lado, la \emph{cadena lateral} es una cadena de átomos la cual es diferente para cada aminoácido, por lo que las características de cada cadena \emph{varían entre aminoácidos}.

\marginpar{
\begin{minipage}{\marginparwidth}
\begin{figure}[H]
\centering
    
    \begin{subfigure}{\marginparwidth}
        \centering
        {\chemfig{H-[7]{\color{blue}N}(-[5]H)-C(-[2]{\color{green}R})-C(=[1]{\color{red}O})-[7]{\color{red}OH}} \par}
        \caption{Estructura general de un aminoácido}
        \label{fig:estructura_aminoácido_general}
    \end{subfigure}
    
    \vspace{2em}
    
    \begin{subfigure}{\marginparwidth}
        \centering
        {\chemfig{H-[7]{\color{blue}N}^{+}(-[5]H)(-[4]H)-C(-[2]{\color{green}R})-C(=[1]{\color{red}O})-[7]{\color{red}OH}} \par}
        \caption{Forma catiónica (+) de un aminoácido, producto de un ambiente ácido}
        \label{fig:estructura_aminoácido_acido}
    \end{subfigure}
    
    \vspace{2em}
    
    \begin{subfigure}{\marginparwidth}
        \centering
        {\chemfig{H-[7]{\color{blue}N}(-[5]H)-C(-[2]{\color{green}R})-C(=[1]{\color{red}O})-[7]{\color{red}O}^{-}} \par}
        \caption{Forma aniónica (-) de un aminoácido, producto de un ambiente básico}
        \label{fig:estructura_aminoácido_básico}
    \end{subfigure}
    
    \vspace{1em}
    \rule{\marginparwidth}{0.2pt}
    \caption{Estructura de un aminoácido y sus variaciones}
\end{figure}
\end{minipage}
}

Así pues, las proteínas son una \emph{cadena de aminoácidos} unidos entre sí por enlaces peptídicos. Sin embargo, esta cadena no se mantiene com una simple linea recta, sino que \emph{se pliega sobre sí misma}, como si se tratara de unos audífonos con cable o unas agujetas. Estos pliegues no ocurren al azar, sino que son resultado de las \emph{interacciones de atracción o repulsión} entre los propios aminoácidos de la cadena, interacciones las cuales provienen de los átomos presentes en la cadena lateral de cada aminoácido. Un ejemplo característico es el de la cisteína, un aminoácido que cuenta con un grupo tiol (\ce{-SH}) en su cadena lateral y que puede formar enlaces disulfuro (\ce{R-S-S-R}) con otras cisteinas presentes en la cadena de aminoácidos. Estos enlaces, como cualquier otro tipo de enlace, están sujetos a las condiciones del medio, tales como su pH, temperatura, etcétera.

De todo lo anterior, lo \emph{esencial} a tener en cuenta al estudiar enzimas es que su estructura y su funcionamiento se derivan de los \emph{aminoácidos que la componen} y la \emph{integridad de sus enlaces}. El conocer su composición de aminoácidos es útil para predecir su plegamiento o identificar los aminoácidos clave involucrados en las reacciones, mientras que cuidar la integridad de sus enlaces es de suma importancia para que las enzimas no sufran ``daños'' que puedan reducir su efectividad para acelerar reacciones.

\section{Morfología de las enzimas}

Como se dijo antes, las enzimas con catalizadores biológicos que aceleran la velocidad de una reacción. Un aspecto clave de esto es que las enzimas cuentan con una \emph{alta especificidad}, o sea, que cada enzima actua de manera óptima sobre un \emph{sustrato específico} y lleva a cabo una \emph{reacción específica}. Si bien algunas enzimas son más flexibles y pueden actuar sobre varios sustratos, generalmente funcionan de manera óptima sobre uno en específico. Cual sea el caso, esta especificidad de las enzimas se origina en el \emph{centro de fijación del sustrato}, el cual es el lugar de la enzima donde el sustrato se va a unir. En este sitio se encuentran varios aminoácidos cuyas cadenas laterales están libres, y son estas cadenas las que forman un enlace con el sustrato y lo fijan a la enzima para que la catálisis pueda llevarse a cabo. Cuáles aminoácidos están presentes en tal sitio es lo que determina a qué sustrato se puede o no unir la enzima.

Una convención común a tener en cuenta antes de avanzar es que esas \emph{cadenas laterales libres} de los aminoácidos que componen a las enzimas (y a las proteínas en general) son más comunmente llamadas \emph{residuos}.

El lugar donde la catálisis se lleva a cabo se conoce como \emph{sitio activo}, y tal catálisis es llevada a cabo de igual manera por ciertos residuos presentes ahí. A estos residuos se les da el nombre de \emph{residuos catalíticos}.

\begin{figure}[H]
    \centering
    \resizebox{\linewidth}{!}{\begin{tikzpicture}
\coordinate (sitio alosterico) at (1.75,1.5);
\coordinate (centro de fijacion 1) at (8,7);
\coordinate (centro de fijacion 2) at (8,5.5);
\coordinate (sitio activo) at (10,4.5);
\coordinate (centro de fijacion 3) at (12,5.5);
\coordinate (centro de fijacion 4) at (12,7);

\draw[line width=1pt] (0,0) 
    to[out=90,in=260] (sitio alosterico)
    to[out=80,in=260] (1,3)
    to[out=90,in=180] (centro de fijacion 1)
    to[out=0,in=90] (centro de fijacion 2)
    to[out=270,in=180] (sitio activo) 
    to[out=0,in=270] (centro de fijacion 3)
    to[out=90,in=180] (centro de fijacion 4)
    to[out=0,in=90] (20,0)
    to cycle;

\node[draw, anchor=north, inner sep=0.75em, align=center, font=\sffamily\large, yshift=-10pt] (sitio activo etiqueta) at (sitio activo.south) {Sitio activo};
\node[draw, anchor=west,  inner sep=0.75em, align=center, text width=4.25em, font=\sffamily\small, xshift=1em] (sitio alosterico etiqueta) at (sitio alosterico.west) {Sitio alostérico};
\node[draw, dashed, rounded corners=4pt, inner sep=0.5em, align=center, text width=6em, fill=white, font=\sffamily\footnotesize] (cambio conformacional) at (6.5,1.5) {Cambio conformacional};

\draw[dashed] (sitio alosterico etiqueta.east) to[out=0,in=180] (cambio conformacional.west);
\draw[-Stealth, dashed] (cambio conformacional.east) to[out=0,in=-90] (sitio activo etiqueta.south);

\node[draw, circle, font=\sffamily, xshift=-3cm] (ligando) at (sitio alosterico) {Ligando};
\draw[-Stealth, dashed] (ligando.east) --++ (1.5,0);

\node[draw, fill=white, circle, line width=0.5pt, node font=\sffamily, xshift=-1cm, yshift=2cm] (sustrato) at (centro de fijacion 1) {Sustrato};
\node[draw, fill=white, circle, line width=0.5pt, node font=\sffamily, xshift=1cm,  yshift=2cm] (producto) at (centro de fijacion 4) {Producto};

\node[draw, align=center, text width=4.5em, font=\sffamily\footnotesize] (residuos cataliticos) at (10,6.5) {Residuos catalíticos};

\node[anchor=west, scale=0.175] (residuo1) at (centro de fijacion 2) {\chemfig{-[:-30]-[:30]*6(-=-(-OH)=-=)}};
\node[anchor=south, scale=0.175] (residuo2) at (sitio activo) {\chemfig{-[:90](-[3]CH3)-[1]CH3}};
\node[anchor=east, scale=0.175] (residuo3) at (centro de fijacion 3) {\chemfig{-[:210]-[:150]-[:210](=[:165]O)-[:255]OH}};

\draw[<-, densely dashed, line width=0.2pt] ([xshift=0.75em]residuo1.east)  to[out=0,in=270]   (residuos cataliticos);
\draw[<-, densely dashed, line width=0.2pt] ([yshift=0.75em]residuo2.north) to[out=90,in=270]  (residuos cataliticos);
\draw[<-, densely dashed, line width=0.2pt] ([xshift=-0.75em]residuo3.west) to[out=180,in=270] (residuos cataliticos);

\draw[-Stealth, dashed] (sustrato.east) to[out=0,in=135] (residuos cataliticos.north west);
\draw[-Stealth, dashed] (residuos cataliticos.north east) to[out=45,in=180] (producto.west);
\end{tikzpicture}}
    \rule{\linewidth}{0.2pt}
    \caption{Esquema general de una enzima y su funcionamiento}
    \label{fig:esquema_de_una_enzima}
\end{figure}

Un matiz importante a mencionar aquí es que el centro de fijación y el sitio activo pueden o no ser el mismo. Cuando son el mismo el proceso es sencillo: el sustrato se une a los residuos del centro de fijación y enseguida esos mismos residuos son los residuos catalíticos que llevarán a cabo la reacción. Cuando el centro de fijación y el sitio activo son diferentes, la enzima tiene que realizar un cambio en su estructura (tambien llamado cambio conformacional) para ``acercar'' el sitio activo y sus residuos catalíticos al sustrato ya fijado.

La última región clave a conocer es el \emph{sitio alostérico}, el cual es una región que solo está presente en algunas enzimas, y la unión de ciertas moléculas en ella induce una amplia gama de efectos, como cambios estructurales, aumento o disminución en el rendimiento de la enzima, completa inhibición de la enzima, etcétera.

\section{Cofactores}

Algunas enzimas necesitan de moléculas extra para funcionar de manera correcta. Estas moléculas son conocidas como \emph{cofactores}. Tales cofactores pueden ser tanto moléculas orgánicas como inorgánicas. A estas enzimas que necesitan de un cofactor para funcionar se les conoce como \emph{apoenzima}, y una vez que tal apoenzima se une con su correspondiente cofactor, entonces toma el nombre de \emph{halo enzima}.